\section{浮动体}
\subsection{图片}
注意：分辨率较高的图片会增加编译时长。
图片格式一般最好用pdf,eps等矢量图，放大不会失真,pdf图片兼容性好，图\ref{fig:pic}。
\begin{figure}[!ht]
    \centering
    \begin{subfigure}[b]{0.32\textwidth}
        \includegraphics[width=\textwidth]{figures/float_exp_ai}
        \caption{PDF}
        \end{subfigure}
    \begin{subfigure}[b]{0.32\textwidth}
        \includegraphics[width=\textwidth]{figures/float_exp_meizi}
        \caption{EPS}
    \end{subfigure}
    \begin{subfigure}[b]{0.32\textwidth}
        \includegraphics[width=\textwidth]{figures/float_exp_mihoyo}
        \caption{JPG}
    \end{subfigure}
    \caption{网络素材}
    \label{fig:pic}
\end{figure}

\subsection{表格}

长表格一般手动填写起来比较费劲，推荐使用工具，如\url{https://tablesgenerator.com/},实在不行用overleaf 自带的也行，长表格如果实在需要断页，一般来说需要标注续表。


\begin{table}[!ht]
\centering % 表格居中
\caption{三线表格} % 表格标题
\label{tab:example} % 用于引用的表格标签
\begin{tabular}{lcr} % lcr表示三列分别左对齐、居中、右对齐
\toprule % 顶部线
\emphasize{左}对齐 & 居中 & \emphasize{右}对齐\\
\midrule % 中间线
$a_1$   & $\int$        &           $\aleph$         \\
$a_1,b_2$   &  $\sum\sum$       &        $ \blacktriangleright \blacktriangleleft$    \\
$a_1,b_2,c_3$    & $\prod\prod\prod$    &   $ \bigstar \circledS  \circlearrowright$   \\
\bottomrule % 底部线
\end{tabular}
\end{table}

\subsection{算法}
算法目前推荐使用\emphasize{algorithm2e}宏包，代码环境一般用
\emphasize{listings}宏包并使用\verb|\lstset|设置风格。

\begin{algorithm}[H]
\DontPrintSemicolon % 不显示分号
\KwData{输入数据}
\KwResult{输出结果}

\SetKwFunction{FMain}{Main}
\SetKwFunction{FSum}{Sum}

\SetKwProg{Fn}{Function}{:}{}
\Fn{\FMain{}}{
  \KwData{这里可以描述函数的输入}
  \KwResult{这里可以描述函数的输出}
  \FSum{$a,b$}\;
  \KwRet{$a+b$}\;
}

\SetKwProg{Fn}{Function}{:}{\KwRet}
\Fn{\FSum{$a, b$}}{
  \KwData{这里可以描述函数的输入}
  \KwResult{这里可以描述函数的输出}
  $sum \gets a + b$\;
  \KwRet{$sum$}\;
}
\caption{示例算法}
\end{algorithm}


\subsection{代码}

\begin{lstlisting}[style=Matlab,caption={MATALB}]
%双曲抛物面
clear;
x = linspace(-10, 10, 100);
y = linspace(-10, 10, 100);
[X, Y] = meshgrid(x, y);
Z = (X.^2)/3 - (Y.^2)/5;
surf(X, Y, Z);  
\end{lstlisting}

\subsection{其他浮动体}
据说支持JavaScript的PDF阅读器，通过\emphasize{media9}宏包，能实现视频或动画的播放。当然利用TikZ(在本文第5节介绍)一帧一帧的放也能实现动画效果，与beamer播放类似。
\clearpage
